\begin{code}[hide]
module slides.popl.sections.example where

open import Data.Nat using (ℕ ; suc)
open import Data.Product using (_×_ ; _,_ ; proj₁ ; proj₂)
open import Data.Unit using (⊤)
open import Level using (0ℓ)
open import Relation.Binary using (Rel)

open import Notations using (₁₊ ; ₂₊ ; ₃₊)
open import Word.Base using (Word ; WRel ; [_]ʷ ; ε ; _•_ ; wmap)
open import Presentation.Definitions using (_IsPresentationOf_)
open import Examples.Groups.Symmetric.Tight.Semantics using (Permutation′-group)
import Examples.Groups.Symmetric.Theorems as SymThm

postulate
  proof-elided : ∀ {ℓ} {A : Set ℓ} → A

private variable
  n : ℕ
\end{code}
%% The running example: permutation circuits (paper §2).
\section{A running example: swap circuits}

\begin{frame}{Circuits are words of wire-indexed generators}
\begin{columns}[c]
\begin{column}{0.58\textwidth}
A gate set is a family of gates by arity. Here: one gate, the swap.
\begin{code}
data Gate : ℕ → Set where
  σ-gate : Gate 2
data Gen : ℕ → Set where
  gate₁ : Gate 1 → Gen (₁₊ n)  -- on the bottom wire
  gate₂ : Gate 2 → Gen (₂₊ n)  -- on the bottom two wires
  _↥   : Gen n → Gen (₁₊ n)    -- one wire higher
\end{code}
\begin{code}[hide]
Circuit : ℕ → Set
Circuit n = Word (Gen n)

CRel : ℕ → Set₁
CRel n = Rel (Circuit n) 0ℓ

_↑ : Circuit n → Circuit (₁₊ n)
_↑ = wmap _↥

pattern σ-gen = gate₂ σ-gate

σ : Circuit (₂₊ n)
σ = [ σ-gen ]ʷ

infix 4 _SRel,_===_ _VRel,_===_ _≈_
infixr 10 σ•_

private variable
  w v : Circuit n

postulate
  _≈_ : Circuit n → Circuit n → Set
\end{code}
\vspace{4pt}
$\aid{Circuit}\;n = \aid{Word}\;(\aid{Gen}\;n)$: a single generator, the empty
circuit $\varepsilon$, or a composition $c\bullet c'$. One constructor of
\aid{Gen} serves \alert{every width} from two up.
\end{column}
\begin{column}{0.38\textwidth}
\centering
\renewcommand{\arraystretch}{1}
\begin{tabular}{@{}r@{\quad}c@{}}
$\sigma$ & \swaps{3}{0} \\
\noalign{\vskip 14pt}
$\sigma\,{\uparrow}$ & \swaps{3}{1} \\
\noalign{\vskip 14pt}
$\sigma\bullet\sigma\,{\uparrow}\bullet\sigma$ & \swaps{3}{0,1,0}
\end{tabular}

\vspace{10pt}
{\scriptsize wire 0 at the bottom; time flows left to right;\\
$c\,{\uparrow} = \aid{wmap}\;\aid{\_↥}\;c$}
\end{column}
\end{columns}
\end{frame}

\begin{frame}{Two axioms --- and structure for free}
\begin{columns}[c]
\begin{column}{0.64\textwidth}
\begin{code}
data _SRel,_===_ : (n : ℕ) → WRel (Gen n) where
  order       : (₂₊ n) SRel, σ • σ === ε
  yang-baxter : (₃₊ n) SRel, σ • σ ↑ • σ === σ ↑ • σ • σ ↑
\end{code}
\vspace{2pt}
{\small Stated once, at the bottom of the wires, for every width. The structural
rules every circuit theory needs come from one generic \aid{Lift-Relation}:}
\renewcommand{\AgdaCodeStyle}{\scriptsize}
\begin{code}
data _VRel,_===_ : (n : ℕ) → CRel n where
  srel  : n SRel, w === v → n VRel, w === v
  cong↑ : n VRel, w === v → (₁₊ n) VRel, w ↑ === v ↑
  comm₁ : (h : Gate 1) (g : Gen n) →
          (₁₊ n) VRel, [ g ↥ ]ʷ • [ gate₁ h ]ʷ === [ gate₁ h ]ʷ • [ g ↥ ]ʷ
  comm₂ : (h : Gate 2) (g : Gen n) →
          (₂₊ n) VRel, [ g ↥ ↥ ]ʷ • [ gate₂ h ]ʷ === [ gate₂ h ]ʷ • [ g ↥ ↥ ]ʷ
\end{code}
\end{column}
\begin{column}{0.33\textwidth}
\centering
\aid{order}\\[3pt]
\fig[0.36]{sn-order}\\[12pt]
\aid{yang-baxter}\\[3pt]
\fig[0.36]{sn-yb}\\[14pt]
{\small shift and disjoint wires commute: the two traces a PROP leaves at
fixed width.}
\end{column}
\end{columns}
\end{frame}

\begin{frame}{Cosets: which wire ends up at the bottom?}
\begin{columns}[c]
\begin{column}{0.47\textwidth}
$S_n$ sits in $S_{n+1}$ as the circuits that leave the bottom wire
alone: the circuits $c\,{\uparrow}$. A coset is determined by \alert{which wire
comes down}, so there are $n+1$ of them, each represented by a
\alert{staircase}.
\begin{code}
data C : ℕ → Set where
  ε   : C n
  σ•_ : C n → C (₁₊ n)
[_]ᶜ : C n → Circuit (₁₊ n)
[ ε ]ᶜ     = ε
[ σ• c ]ᶜ  = σ • ([ c ]ᶜ ↑)
\end{code}
\end{column}
\begin{column}{0.49\textwidth}
\centering
\begin{tabular}{@{}l@{\;\;}ccc@{}}
  $\aid{C}\;2$: & \fig[0.42]{c2-0} & \fig[0.42]{c2-1} & \fig[0.42]{c2-2}\\
  & $\varepsilon$ & $\sigma{\bullet}\,\varepsilon$ & $\sigma{\bullet}\,\sigma{\bullet}\,\varepsilon$\\[10pt]
  $\aid{C}\;1$: & \fig[0.42]{c1-0} & \fig[0.42]{c1-1} & \\
  & $\varepsilon$ & $\sigma{\bullet}\,\varepsilon$ & \\[10pt]
  $\aid{C}\;0$: & \fig[0.42]{c0-0} & & \\
  & $\varepsilon$ & &
\end{tabular}

\vspace{10pt}
\begin{idea}
\small One digit per wire: $|S_{n+1}| = (n{+}1)\cdot n\cdots 2$ --- the
\alert{factorial number system}.
\end{idea}
\end{column}
\end{columns}
\end{frame}

\begin{frame}{The coset table: push one generator past a staircase}
\begin{columns}[c]
\begin{column}{0.66\textwidth}
\renewcommand{\AgdaCodeStyle}{\footnotesize}
\begin{code}
ract : C (₁₊ n) → Gen (₂₊ n) → Circuit (₁₊ n) × C (₁₊ n)
ract ε σ-gen = ε , σ• ε             -- grows
ract (σ• ε) σ-gen = ε , ε           -- cancels (σ² = ε)
ract (σ• σ• c) σ-gen = σ , σ• σ• c  -- passes (Yang-Baxter)
ract ε (g ↥) = [ g ]ʷ , ε           -- commutes past
ract (σ• c) (g ↥) = proj₁ (ract c g) ↑ , σ• (proj₂ (ract c g))
\end{code}
\vspace{2pt}
{\small Each clause is certified by its \alert{soundness law}:}
\begin{code}
ract-sound : ∀ c b → let (b' , c') = ract c b in
             [ c ]ᶜ • [ b ]ʷ ≈ b' ↑ • [ c' ]ᶜ
\end{code}
\begin{code}[hide]
ract-sound = proof-elided
\end{code}
\end{column}
\begin{column}{0.31\textwidth}
\centering
\fig[0.36]{ract-l}\;$\approx$\;\fig[0.36]{ract-r}\\[4pt]
{\scriptsize staircase, then generator\\$\approx$ residual one wire up,
then staircase}
\end{column}
\end{columns}
\vspace{4pt}
\centering
\begin{tabular}{@{}c@{\hspace{7mm}}c@{\hspace{7mm}}c@{\hspace{7mm}}c@{}}
\fig[0.27]{push4-l}\,$\to$\,\fig[0.27]{push4-r} &
\fig[0.27]{push3-l}\,$\to$\,\fig[0.27]{push3-r} &
\fig[0.27]{push2-l}\,$\to$\,\fig[0.27]{push2-r} &
\fig[0.27]{push1-l}\,$\to$\,\fig[0.27]{push1-r}\\[2pt]
\scriptsize grows & \scriptsize cancels & \scriptsize passes (Yang--Baxter) & \scriptsize commutes past
\end{tabular}
\end{frame}

\begin{frame}{A tower of tables, and the theorem}
\begin{columns}[c]
\begin{column}{0.50\textwidth}
Threading the table through a circuit gives one level of normalization;
iterating it, one level per wire, is the \alert{coset tower}:
\begin{code}
NF : ℕ → Set
NF 0       = ⊤
NF (suc k) = NF k × C k   -- ⊤ × C 0 × C 1 × ⋯ × C k
\end{code}
\vspace{2pt}
\centering
\fig[0.38]{fig3_staircase}\quad{\scriptsize a ``staircase of staircases''}
\end{column}
\begin{column}{0.47\textwidth}
\small\aid{nf} runs the tables; \aid{inv-nf} rebuilds a circuit from its digits.
If $\llbracket w\rrbracket = \llbracket v\rrbracket$:
\begin{center}
\begin{tikzpicture}[font=\small]
  \node (w) at (0,1.2) {$w$};
  \node (v) at (4.0,1.2) {$v$};
  \node (nw) at (0,0) {\aid{inv-nf (nf w)}};
  \node (nv) at (4.0,0) {\aid{inv-nf (nf v)}};
  \draw[double,thick,prim] (w) -- node[left,font=\scriptsize]{$\approx$} (nw);
  \draw[double,thick,prim] (v) -- node[right,font=\scriptsize]{$\approx$} (nv);
  \draw[double,thick,ok] (nw) -- (nv);
  \node[font=\scriptsize,text=ok] at (2.0,-0.45) {$\equiv$ by uniqueness};
  \node[font=\scriptsize,text=prim] at (2.0,1.2) {normalize both};
\end{tikzpicture}
\end{center}
--- a derivation built from \aid{order}, \aid{yang-baxter} and the structural
rules only. Even the \emph{endofunction} semantics is complete.
\end{column}
\end{columns}
\vspace{6pt}
\begin{code}
symmetric-presentation : ∀ n → (n VRel,_===_) IsPresentationOf (Permutation′-group n)
\end{code}
\begin{code}[hide]
symmetric-presentation n = SymThm.Tight.presentation n
\end{code}
\end{frame}
